% Generated draft: requires mathematical and source review.
\documentclass[11pt]{article}
\usepackage[margin=25mm]{geometry}
\usepackage{fontspec}
\setmainfont{Latin Modern Roman}
\usepackage{amsmath,amssymb,mathtools,unicode-math}
\setmathfont{Latin Modern Math}
\usepackage{enumitem,needspace,xurl,hyperref}
\setlength{\parindent}{0pt}
\setlength{\parskip}{5pt}
\setlength{\emergencystretch}{3em}
\newcommand{\N}{\mathbb{N}}
\newcommand{\Z}{\mathbb{Z}}
\newcommand{\Q}{\mathbb{Q}}
\newcommand{\R}{\mathbb{R}}
\newcommand{\C}{\mathbb{C}}
\newcommand{\F}{\mathbb{F}}
\newcommand{\A}{\mathbb{A}}
\newcommand{\PP}{\mathbb{P}}
\newcommand{\E}{\mathbb{E}}
\newcommand{\eps}{\varepsilon}
\newcommand{\dd}{\,\mathrm{d}}
\begin{document}
\section*{UnsolvedMath ID: 7000026}
\textbf{Source problem code:} catalogue entry
\subsection{Definitions and mathematical statement}
In the usual complete-surface convention, let \(F:\Sigma^2\to\R^3\) be a smooth immersion of a connected surface without boundary, complete for the induced metric \(F^*g_{\mathrm{Eucl}}\). Let \(k_1,k_2\) denote its principal curvatures, and \(K=k_1k_2\) its Gaussian curvature. Absolute values do not depend on the choice of unit normal.

The question asks whether there exist such a surface and a constant \(c>0\) with
\[
 K(p)<0\quad\text{and}\quad |k_i(p)|\geq c
 \qquad(p\in\Sigma,\ i=1,2).
\]
The bound is global and uniform for both principal curvatures; mere nonvanishing at each point is weaker. The ambient embedding requirement and precise differentiability class are not specified by the extracted phrase ``surfaces in Euclidean 3-space''; the immersion formulation makes the standard broader setting explicit for review.

The registry labels the answer negative and attributes it to Efimov's theorem (1963), reported there as excluding complete surfaces with \(K\leq-c_0<0\). The displayed principal curvature bounds imply \(K\leq-c^2\), so that reported theorem would address this formulation if its regularity and completeness hypotheses match. The converse implication is not asserted: a lower bound on \(|k_1k_2|\) alone does not give separate positive lower bounds on both factors. Thus the summary's claimed equivalence of these two uniform conditions is not adopted.
\subsection{Short English statement}
Can a complete negatively curved surface in \(\R^3\) have both principal curvature magnitudes uniformly bounded below by one positive constant? The registry reports a negative answer via Efimov; its asserted equivalence of separate principal bounds with a Gaussian bound needs correction.
\subsection{Sources}
\begin{itemize}
\item Local UnsolvedMath registry, permanent ID 7000026, source code catalogue entry; attributed to Mohammad Ghomi, proposed year 2019; update 2026-08-06; dated status solved.
\item Ghomi, \emph{Open Problems in Geometry of Curves and Surfaces} (2019), Problem 7.1, PDF page 15: \url{https://people.math.gatech.edu/~ghomi/Papers/op.pdf}. Registry access date 2026-07-29; normalized PDF extraction; public author-hosted PDF; current status and reuse terms require review.
\item The supplied summary cites Efimov's theorem (1963), without an exact paper citation, as excluding complete surfaces in \(\R^3\) with Gaussian curvature bounded above by a negative constant. The attribution and exact hypotheses have not been independently verified here.
\end{itemize}
\end{document}
